
% resumo em português OBRIGATÓRIO
\setlength{\absparsep}{18pt} % ajusta o espaçamento dos parágrafos do resumo
\begin{resumo}
Esta proposta de desenvolvimento de atividade de extensão propõe o desenvolvimento de um protótipo de bracelete inteligente baseado na plataforma ESP32, concebido como tecnologia assistiva de baixo custo para auxiliar pessoas com baixa visão na percepção espacial e na orientação diante de obstáculos no ambiente. A locomoção independente de indivíduos com deficiência visual enfrenta frequentes desafios diante de barreiras físicas e obstáculos suspensos que não são detectados pela bengala convencional. Diante desse cenário, o dispositivo vestível a ser desenvolvido integrará sensores de proximidade e um atuador vibrotátil ao microcontrolador ESP32, visando converter leituras de distância em alertas táteis intuitivos e não invasivos no punho do usuário. O objetivo geral consiste em desenvolver e testar o protótipo, avaliando sua eficácia na detecção de obstáculos e na facilidade de interpretação dos estímulos gerados. A metodologia planejada será estruturada em etapas consecutivas, contemplando o levantamento de requisitos com base nas demandas de acessibilidade, a pesquisa e seleção dos componentes eletrônicos, o projeto e montagem da estrutura física do bracelete, a implementação do software embarcado para processamento dos sinais e a realização de testes controlados de funcionamento. Pretende-se avaliar o desempenho do equipamento, mapear suas limitações técnicas e propor melhorias futuras, promovendo a integração entre a formação acadêmica em Ciência da Computação e as demandas sociais de inclusão e acessibilidade.
 
 \textbf{Palavras-chave}: Tecnologia assistiva; ESP32; Internet das Coisas; Baixa visão; Sistemas embarcados.
\end{resumo}

% resumo em inglês OBRIGATÓRIO
\begin{resumo}[Abstract]
 \begin{otherlanguage*}{english}
This extension activity proposal puts forward the development of a smart wristband prototype based on the ESP32 platform, conceived as a low-cost assistive technology to support individuals with low vision in spatial perception and navigation around obstacles in their environment. Independent mobility for visually impaired individuals frequently faces challenges posed by physical barriers and overhead obstacles that remain undetected by conventional white canes. In light of this scenario, the wearable device to be developed will integrate proximity sensors and a vibrotactile actuator with the ESP32 microcontroller, aiming to convert distance measurements into intuitive and non-invasive haptic alerts on the user's wrist. The general objective is to develop and test the prototype, assessing its efficacy in obstacle detection and the ease of interpreting the generated stimuli. The planned methodology will be structured into sequential stages, encompassing requirements elicitation based on accessibility needs, research and selection of electronic components, design and physical assembly of the wristband structure, implementation of embedded software for signal processing, and the execution of controlled functional tests. It is intended to evaluate device performance, identify technical limitations, and propose future improvements, fostering the integration between academic education in Computer Science and social demands for inclusion and accessibility.

   \vspace{\onelineskip}
 
   \noindent 
   \textbf{Keywords}: Assistive technology; ESP32; Internet of Things; Low vision; Embedded systems.
 \end{otherlanguage*}
\end{resumo}
